\section{45维构造所用母格的最小范数}
\label{app:ambient}


母格的具体身份由如下构造固定。令 \(C\) 为扩展三元二次剩余 \([48,24,15]\)
码 \cite{nebe}， \(M=\{z\in\mathbb Z^{48}:z\bmod3\in C\}\)，
\(L_0=\{z/\sqrt3:z\in M,\ z^2\equiv0\pmod6\}\)。
由所给整数基求出特征向量分子 \(\chi\in M\)，所有坐标均为奇数，满足
\(\chi^2=264\) 及 \(\chi\cdot z\equiv z^2\pmod6\)。取

\[
\Lambda=L_0\cup\bigl(\chi/(2\sqrt3)+L_0\bigr).
\]

特征同余与 \(\chi^2/12=22\)
保证整性和偶性；两个指数二的构造保证余体积为一。程序从这两个陪集重建整数基
\(B\)，检查 \(BB^{\mathsf T}=12G\)，其中 \(G\) 正是298个见证所用的 Gram
矩阵。

在 \(L_0\) 中，范数2或4的向量要求分子 \(z^2\le12\)。码的最小重量15迫使
\(z\bmod3=0\)，此时非零偶范数至少为6。在另一陪集中，分子 \(y=\chi+2z\)
的48个坐标均为奇数，所以 \(y^2/12\ge4\)。若等于4，则每个坐标均为
\(\pm1\)。

令 \(s_i\in\{\pm1\}\) 为 \(\chi_i\bmod3\) 的代表；有限检查给出
\(s\in C\)、\(\chi\cdot s=66\) 及 \(\chi_i s_i\equiv1\pmod6\)。按 \(s\)
变换坐标后，码包含全一字。其二符号重量枚举式为 \(1+94t^{24}+t^{48}\)
\cite[equation (20)]{munemasa-tamura}，所以二进制取值码字的重量只能为0、24、48。写
\(y=s(1-2b)\)，得到二进制码字 \(b\) 及重量 \(w\)，但

\[
\chi\cdot z
=\frac{\chi\cdot s-\chi^2}{2}-\sum_i\chi_i s_i b_i
\equiv-99-w\equiv3\pmod6,
\]

与 \(z/\sqrt3\in L_0\) 矛盾。偶性排除其余小于6的正范数，而
\(\sqrt3(e_1+e_2)\in L_0\) 的平方范数恰为6。于是上述 theta
级数和设计矩均适用于这一个具体母格。
